\subsection*{Гипотеза простого равномерного хеширования.}

Проблема хеш-таблицы на цепочках - неравномерная длинна цепочки из-за неравномерности возникновения коллизии. $\Rightarrow$ Хочется минимизировать максимальную длину цепочки, то есть добиться того чтобы функция равномерно распределяла ключи по бакетам.

Пусть $\mathcal{H} = \{h:U \rightarrow \{0, 1, \dots, k - 1\} \}$ - множество хеш-функций.

\Note $U \subset \mathbb{N}$ и $|U| < \infty$, то есть $U = \{0, 1, \dots, n - 1\}$.

Определим случайную функцию как то, что образ каждого ключа - это случайный элемент из множества $\{0, 1, \dots, k - 1\}$.

\Note Неприменимо на практике, т.к. требует $\mathcal{O}(|U|\log|h(U)|)$ памяти на хранение такого отображения.

\Def Модель, описанная выше, называется \textit{простое равномерное хеширование} (\textit{simple uniform hashing}).

\subsection*{Математическое ожидание длины цепочки.}

Величина $L_q$ - длинна цепочки по ключу $q$.\newline
Пусть построена таблица для ключей $k_1, \dots, k_n$, тогда 

\begin{center}
{\large $L_q = \sum\limits_{i = 1}^{n} I(h(q) = h(k_i))$}
\end{center}

Посчитаем матожидание длины цепочки.

\begin{center}
{\large $\mathbb{E} L_q = \mathbb{E}\sum\limits_{i = 1}^{n} I(h(q) = h(k_i)) = \sum\limits_{i = 1}^{n} P(h(q) = h(k_i))$}
\end{center}

Посчитаем вероятность.

\begin{center}
    {\large $ P(h(x) = h(y)) =
        \begin{cases}
        1, \ x = y\\
        \frac{1}{k}, \ x \neq y 
    \end{cases} $}
\end{center}

С учетом того что все ключи различны, получаем, что

\begin{center}
{\large $\mathbb{E} L_q = \mathbb{E} \sum\limits_{i = 1}^{n} I(h(q) = h(k_i)) = \sum\limits_{i = 1}^{n} P(h(q) = h(k_i)) \leq 1 + \frac{n}{k}$}
\end{center}

\Def Коэффициентом загруженности \textit{(load factor)} называют величину $\alpha = \frac{n}{k}$.

\Note Из оценок выше можно задать константу $C$ такую, что всегда $\alpha < C$. Тогда в среднем сложность операций составит $\mathcal{O}(1)$.

\pagebreak